\documentclass[11pt,letterpaper]{article}
\usepackage[margin=1in]{geometry}
\usepackage[T1]{fontenc}
\usepackage{lmodern,microtype,booktabs,graphicx,amsmath}
\usepackage[round,authoryear]{natbib}
\usepackage[font=small,labelfont=bf]{caption}
\usepackage[hidelinks]{hyperref}
\usepackage{xurl,setspace}
\setstretch{1.15}
\title{Big, If True}
\author{Andrew Gelman \and Gaurav Sood}
\date{Working paper, \today}
\hypersetup{pdfauthor={Andrew Gelman and Gaurav Sood},pdftitle={Big, If True}}
\begin{document}
\maketitle
\begin{abstract}
\noindent To be written.
\end{abstract}

\section{Introduction}
To be written. Low-powered studies that clear a significance threshold
overstate effect sizes \citep{gelman2014}.

\section{Cases}

\subsection{Emotional and rational leaflets in Allentown}
\citet{westen2007} reports that in a 1930s field study in Allentown,
Pennsylvania, Socialist votes rose by 50 percent in a ward that received an
``emotional'' leaflet, by about a third in a ward that received a ``rational''
leaflet, and by about a quarter in a ward that received none. The source is
\citet{hartmann1936}.

\bibliographystyle{plainnat}
\bibliography{references}
\end{document}
